% GENERATED FILE. Edit canonical lessons, then run npm run generate:books.
% canonical-source-hash: fnv1a64:e5f0d97aa32b8bb8
% canonical-lessons: BN-C453-atmobishvas, BN-C453-monobol, BN-C453-kapurush, BN-C453-rukhe-darano, BN-C453-pichhiye-jaoya, BN-R453-review-446-453

\chapter{Courage and Nerve}
\label{ch:bn-courage-and-nerve}
\hlchaptermodality{453}

\begin{chapteropening}
\textbf{By the end of this chapter:} \emph{I can talk about confidence and about standing up to someone.}

Complete the last lesson of chapter 453: I can talk about confidence and about standing up to someone.
\end{chapteropening}

\section[\texorpdfstring{\bn{আত্মবিশ্বাস} (ātmôbishvās)}{আত্মবিশ্বাস (ātmôbishvās)}]{\bn{আত্মবিশ্বাস} (ātmôbishvās) — self-confidence}

\label{lesson:BN-C453-atmobishvas}



\begin{quote}
Before the new one: say the Bengali for to show off, then the Bengali for to be humiliated.
\end{quote}

\subsection*{You'll want to know: \bn{আত্মবিশ্বাস}}
\textbf{\bn{আত্মবিশ্বাস}} — \emph{ātmôbishvās} — \textquotedblleft{}self-confidence\textquotedblright{}.

\begin{quote}
\bn{গান} \bn{শেখার} \bn{পরে} \bn{মীরার} \bn{আত্মবিশ্বাস} \bn{বাড়ল।}
\end{quote}

\begin{quote}
\emph{gān shekhār pôre mīrār ātmôbishvās bāṛlô.}
\end{quote}

\begin{quote}
After learning to sing, Mira's self-confidence grew.
\end{quote}

\begin{cousinweb}
\textbf{\bn{আত্মবিশ্বাস}} (\emph{ātmôbishvās}) puts Sanskrit \emph{ātman}, \textquotedblleft{}self\textquotedblright{}, in front of \emph{viśvāsa}, \textquotedblleft{}trust\textquotedblright{} (Monier-Williams), the \bn{বিশ্বাস} of \bn{বিশ্বাস} \bn{করা} (taught), to believe: trust in yourself.
\end{cousinweb}

\begin{grammarlens}[title={What you've built}]
One more word for describing what you lived and felt.
\end{grammarlens}

\subsection*{Your turn}
\begin{itemize}
\raggedright
  \item \emph{Say it:} \emph{ātmôbishvās}
  \item \emph{Say it:} \emph{ātmôbishvās}, once more
  \item \emph{Say it:} \emph{gān shekhār pôre mīrār ātmôbishvās bāṛlô.} — After learning to sing, Mira's self-confidence grew.
  \item \emph{Recall:} say the Bengali for to show off, then the Bengali for to be humiliated, then say \emph{ātmôbishvās} again
\end{itemize}

\begin{tcolorbox}[breakable,colback=teal!4,colframe=teal!35!black,title={Before you move on}]
What does \bn{আত্মবিশ্বাস} mean? (\textquotedblleft{}Self-confidence\textquotedblright{}.) Say it once more.
\end{tcolorbox}
\section[\texorpdfstring{\bn{মনোবল} (mônobôl)}{মনোবল (mônobôl)}]{\bn{মনোবল} (mônobôl) — morale; strength of mind}

\label{lesson:BN-C453-monobol}



\begin{quote}
Before the new one: say the Bengali for to be humiliated, then the Bengali for self-confidence.
\end{quote}

\subsection*{You'll want to know: \bn{মনোবল}}
\textbf{\bn{মনোবল}} — \emph{mônobôl} — \textquotedblleft{}morale; strength of mind\textquotedblright{}.

\begin{quote}
\bn{দলটা} \bn{খেলায়} \bn{জিতল} \bn{না}, \bn{কিন্তু} \bn{খেলোয়াড়দের} \bn{মনোবল} \bn{কমল} \bn{না।}
\end{quote}

\begin{quote}
\emph{dôlṭā khelāy jitlô nā, kintu kheloyāṛder mônobôl kômlô nā.}
\end{quote}

\begin{quote}
The team did not win the match, but the players' morale did not drop.
\end{quote}

\begin{cousinweb}
\textbf{\bn{মনোবল}} (\emph{mônobôl}) joins Sanskrit \emph{manas}, \textquotedblleft{}mind\textquotedblright{} (CDIAL 9822), in its compound form \emph{mano-}, and \emph{bala}, \textquotedblleft{}strength\textquotedblright{} (Monier-Williams): strength of mind.
\end{cousinweb}

\begin{grammarlens}[title={What you've built}]
One more word for describing what you lived and felt.
\end{grammarlens}

\subsection*{Your turn}
\begin{itemize}
\raggedright
  \item \emph{Say it:} \emph{mônobôl}
  \item \emph{Say it:} \emph{mônobôl}, once more
  \item \emph{Say it:} \emph{dôlṭā khelāy jitlô nā, kintu kheloyāṛder mônobôl kômlô nā.} — The team did not win the match, but the players' morale did not drop.
  \item \emph{Recall:} say the Bengali for to be humiliated, then the Bengali for self-confidence, then say \emph{mônobôl} again
\end{itemize}

\begin{tcolorbox}[breakable,colback=teal!4,colframe=teal!35!black,title={Before you move on}]
What does \bn{মনোবল} mean? (\textquotedblleft{}Morale; strength of mind\textquotedblright{}.) Say it once more.
\end{tcolorbox}
\section[\texorpdfstring{\bn{কাপুরুষ} (kāpurush)}{কাপুরুষ (kāpurush)}]{\bn{কাপুরুষ} (kāpurush) — a coward}

\label{lesson:BN-C453-kapurush}



\begin{quote}
Before the new one: say the Bengali for self-confidence, then the Bengali for morale.
\end{quote}

\subsection*{You'll want to know: \bn{কাপুরুষ}}
\textbf{\bn{কাপুরুষ}} — \emph{kāpurush} — \textquotedblleft{}a coward\textquotedblright{}.

\begin{quote}
\bn{বিপদ} \bn{দেখে} \bn{পালানো} \bn{কাপুরুষের} \bn{কাজ।}
\end{quote}

\begin{quote}
\emph{bipôd dekhe pālāno kāpurusher kāj.}
\end{quote}

\begin{quote}
Running away at the sight of danger is a coward's act.
\end{quote}

\bn{ভীতু} (taught) is timid; \bn{কাপুরুষ} names the man himself, and it stings.

\begin{cousinweb}
\textbf{\bn{কাপুরুষ}} (\emph{kāpurush}) is a tatsama, Sanskrit \emph{kāpuruṣa}, \textquotedblleft{}a contemptible man, a coward\textquotedblright{} (Monier-Williams).
\end{cousinweb}

\begin{grammarlens}[title={What you've built}]
One more word for describing what you lived and felt.
\end{grammarlens}

\subsection*{Your turn}
\begin{itemize}
\raggedright
  \item \emph{Say it:} \emph{kāpurush}
  \item \emph{Say it:} \emph{kāpurush}, once more
  \item \emph{Say it:} \emph{bipôd dekhe pālāno kāpurusher kāj.} — Running away at the sight of danger is a coward's act.
  \item \emph{Recall:} say the Bengali for self-confidence, then the Bengali for morale, then say \emph{kāpurush} again
\end{itemize}

\begin{tcolorbox}[breakable,colback=teal!4,colframe=teal!35!black,title={Before you move on}]
What does \bn{কাপুরুষ} mean? (\textquotedblleft{}A coward\textquotedblright{}.) Say it once more.
\end{tcolorbox}
\section[\texorpdfstring{\bn{রুখে} \bn{দাঁড়ানো} (rukhe dā̃ṛāno)}{রুখে দাঁড়ানো (rukhe dā̃ṛāno)}]{\bn{রুখে} \bn{দাঁড়ানো} (rukhe dā̃ṛāno) — to stand up to; to resist}

\label{lesson:BN-C453-rukhe-darano}



\begin{quote}
Before the new one: say the Bengali for morale, then the Bengali for a coward.
\end{quote}

\subsection*{You'll want to know: \bn{রুখে} \bn{দাঁড়ানো}}
\textbf{\bn{রুখে} \bn{দাঁড়ানো}} — \emph{rukhe dā̃ṛāno} — \textquotedblleft{}to stand up to; to resist\textquotedblright{}.

\begin{quote}
\bn{মীরা} \bn{একা} \bn{ডাকাতদের} \bn{সামনে} \bn{রুখে} \bn{দাঁড়াল।}
\end{quote}

\begin{quote}
\emph{mīrā ekā ḍākātder sāmne rukhe dā̃ṛālô.}
\end{quote}

\begin{quote}
Mira stood up to the robbers on her own.
\end{quote}

Past (he/she …): \bn{তিনি} \bn{চোরের} \bn{সামনে} \bn{রুখে} \bn{দাঁড়ালেন}, \emph{tini chorer sāmne rukhe dā̃ṛālen}.

\begin{cousinweb}
\textbf{\bn{রুখে} \bn{দাঁড়ানো}} (\emph{rukhe dā̃ṛāno}) is \emph{rukhe}, the -\bn{এ} form of \emph{rokhā}, \textquotedblleft{}to hold back, to resist\textquotedblright{}. CDIAL lists a Bengali \emph{rokā}, to hinder (CDIAL 10827), and a \emph{rukhā}, to be angry (CDIAL 10800); which of them this is, is not settled. \bn{দাঁড়ানো} (taught), to stand, probably belongs with Sanskrit \emph{daṇḍa}, a stick (CDIAL 6128).
\end{cousinweb}

\begin{grammarlens}[title={What you've built}]
One more word for describing what you lived and felt.
\end{grammarlens}

\subsection*{Your turn}
\begin{itemize}
\raggedright
  \item \emph{Say it:} \emph{rukhe dā̃ṛāno}
  \item \emph{Say it:} \emph{rukhe dā̃ṛāno}, once more
  \item \emph{Say it:} \emph{mīrā ekā ḍākātder sāmne rukhe dā̃ṛālô.} — Mira stood up to the robbers on her own.
  \item \emph{Recall:} say the Bengali for morale, then the Bengali for a coward, then say \emph{rukhe dā̃ṛāno} again
\end{itemize}

\begin{tcolorbox}[breakable,colback=teal!4,colframe=teal!35!black,title={Before you move on}]
What does \bn{রুখে} \bn{দাঁড়ানো} mean? (\textquotedblleft{}To stand up to; to resist\textquotedblright{}.) Say it once more.
\end{tcolorbox}
\section[\texorpdfstring{\bn{পিছিয়ে} \bn{যাওয়া} (pichhiye jāoyā)}{পিছিয়ে যাওয়া (pichhiye jāoyā)}]{\bn{পিছিয়ে} \bn{যাওয়া} (pichhiye jāoyā) — to back down; to fall behind}

\label{lesson:BN-C453-pichhiye-jaoya}



\begin{quote}
Before the new one: say the Bengali for a coward, then the Bengali for to stand up to.
\end{quote}

\subsection*{You'll want to know: \bn{পিছিয়ে} \bn{যাওয়া}}
\textbf{\bn{পিছিয়ে} \bn{যাওয়া}} — \emph{pichhiye jāoyā} — \textquotedblleft{}to back down; to fall behind\textquotedblright{}.

\begin{quote}
\bn{মীরা} \bn{ভয়} \bn{পেল}, \bn{কিন্তু} \bn{পিছিয়ে} \bn{গেল} \bn{না।}
\end{quote}

\begin{quote}
\emph{mīrā bhôy pelô, kintu pichhiye gelô nā.}
\end{quote}

\begin{quote}
Mira was afraid, but she did not back down.
\end{quote}

Past (he/she …): \bn{তিনি} \bn{শেষে} \bn{পিছিয়ে} \bn{গেলেন}, \emph{tini sheshe pichhiye gelen}.

\begin{cousinweb}
\textbf{\bn{পিছিয়ে} \bn{যাওয়া}} (\emph{pichhiye jāoyā}) is \emph{pichhiye}, the -\bn{এ} form of \emph{pichhāno}, \textquotedblleft{}to move back\textquotedblright{}, built on the word for the rear that CDIAL lists for Bengali (CDIAL 7990), the \bn{পিছ} of \bn{পিছনে} (taught), behind; \emph{jāoyā} is \textquotedblleft{}to go\textquotedblright{} (CDIAL 10452).
\end{cousinweb}

\begin{grammarlens}[title={What you've built}]
One more word for describing what you lived and felt.
\end{grammarlens}

\subsection*{Your turn}
\begin{itemize}
\raggedright
  \item \emph{Say it:} \emph{pichhiye jāoyā}
  \item \emph{Say it:} \emph{pichhiye jāoyā}, once more
  \item \emph{Say it:} \emph{mīrā bhôy pelô, kintu pichhiye gelô nā.} — Mira was afraid, but she did not back down.
  \item \emph{Recall:} say the Bengali for a coward, then the Bengali for to stand up to, then say \emph{pichhiye jāoyā} again
\end{itemize}

\begin{tcolorbox}[breakable,colback=teal!4,colframe=teal!35!black,title={Before you move on}]
What does \bn{পিছিয়ে} \bn{যাওয়া} mean? (\textquotedblleft{}To back down; to fall behind\textquotedblright{}.) Say it once more.
\end{tcolorbox}
\section[(dialogue)]{Looking back: from Comfort and Relief to Courage and Nerve}

\label{lesson:BN-R453-review-446-453}



\begin{quote}
Say the Bengali for consolation and ease of mind.
\end{quote}

\subsection*{Your turn}
Say these aloud:

\begin{itemize}
\raggedright
  \item consolation, ease of mind, reliance, to breathe a sigh of relief, to cope with
  \item regret, remorse, repentant, to rue, to admit
  \item the pain of separation, a yearning, wistful, anxious with longing, to cry out in despair
  \item affection, tender attachment, cuddling and petting, romantic love, to be devoted to
  \item trifling, mockery, to insult, to taunt, to belittle
  \item revolting, aversion, envious of others' good fortune, to hate, to turn up one's nose
  \item arrogance, conceit, self-respect, to show off, to be humiliated
  \item self-confidence, morale, a coward, to stand up to, to back down
\end{itemize}

\begin{tcolorbox}[breakable,colback=teal!4,colframe=teal!35!black,title={Before you move on}]
Consolation? (\textbf{\bn{সান্ত্বনা}}, \emph{sāntvônā}.) Ease of mind? (\textbf{\bn{স্বস্তি}}, \emph{svôsti}.) Reliance? (\textbf{\bn{ভরসা}}, \emph{bhôrsā}.) To breathe a sigh of relief? (\textbf{\bn{হাঁফ} \bn{ছাড়া}}, \emph{hā̃ph chhāṛā}.) To cope with? (\textbf{\bn{সামলানো}}, \emph{sāmlāno}.) Regret? (\textbf{\bn{আফসোস}}, \emph{āphsos}.) Remorse? (\textbf{\bn{অনুশোচনা}}, \emph{ônushochônā}.) Repentant? (\textbf{\bn{অনুতপ্ত}}, \emph{ônutôptô}.) To rue? (\textbf{\bn{পস্তানো}}, \emph{pôstāno}.) To admit? (\textbf{\bn{স্বীকার} \bn{করা}}, \emph{svīkār kôrā}.) The pain of separation? (\textbf{\bn{বিরহ}}, \emph{birôhô}.) A yearning? (\textbf{\bn{আকাঙ্ক্ষা}}, \emph{ākāngkhā}.) Wistful? (\textbf{\bn{উদাস}}, \emph{udās}.) Anxious with longing? (\textbf{\bn{ব্যাকুল}}, \emph{bækul}.) To cry out in despair? (\textbf{\bn{হাহাকার} \bn{করা}}, \emph{hāhākār kôrā}.) Affection? (\textbf{\bn{স্নেহ}}, \emph{sneh}.) Tender attachment? (\textbf{\bn{মায়া}}, \emph{māyā}.) Cuddling and petting? (\textbf{\bn{আদর}}, \emph{ādôr}.) Romantic love? (\textbf{\bn{প্রেম}}, \emph{prem}.) To be devoted to? (\textbf{\bn{ভক্তি} \bn{করা}}, \emph{bhôkti kôrā}.) Trifling? (\textbf{\bn{তুচ্ছ}}, \emph{tuchchhô}.) Mockery? (\textbf{\bn{উপহাস}}, \emph{upôhās}.) To insult? (\textbf{\bn{অপমান} \bn{করা}}, \emph{ôpômān kôrā}.) To taunt? (\textbf{\bn{খোঁটা} \bn{দেওয়া}}, \emph{khõṭā deoyā}.) To belittle? (\textbf{\bn{হেয়} \bn{করা}}, \emph{heyô kôrā}.) Revolting? (\textbf{\bn{ঘিনঘিনে}}, \emph{ghinghine}.) Aversion? (\textbf{\bn{বিতৃষ্ণা}}, \emph{bitrishnā}.) Envious of others' good fortune? (\textbf{\bn{পরশ্রীকাতর}}, \emph{pôrôshrīkātôr}.) To hate? (\textbf{\bn{ঘৃণা} \bn{করা}}, \emph{ghrinā kôrā}.) To turn up one's nose? (\textbf{\bn{নাক} \bn{সিঁটকানো}}, \emph{nāk sĩṭkāno}.) Arrogance? (\textbf{\bn{অহংকার}}, \emph{ôhôngkār}.) Conceit? (\textbf{\bn{দেমাক}}, \emph{demāk}.) Self-respect? (\textbf{\bn{আত্মসম্মান}}, \emph{ātmôsômmān}.) To show off? (\textbf{\bn{জাহির} \bn{করা}}, \emph{jāhir kôrā}.) To be humiliated? (\textbf{\bn{অপদস্থ} \bn{হওয়া}}, \emph{ôpôdôsthô hôoyā}.) Self-confidence? (\textbf{\bn{আত্মবিশ্বাস}}, \emph{ātmôbishvās}.) Morale? (\textbf{\bn{মনোবল}}, \emph{mônobôl}.) A coward? (\textbf{\bn{কাপুরুষ}}, \emph{kāpurush}.) To stand up to? (\textbf{\bn{রুখে} \bn{দাঁড়ানো}}, \emph{rukhe dā̃ṛāno}.) To back down? (\textbf{\bn{পিছিয়ে} \bn{যাওয়া}}, \emph{pichhiye jāoyā}.)
\end{tcolorbox}
